\begin{figure}[htp!]
  \centering
  \includegraphics[scale=0.22]{Esercizio13.1/Esercizio.jpg} 
\end{figure}
Utilizzando la distribuzione di corpo nero e ricordando la legge di Wien è possibile calcolare l'energia tipica dei fotoni a questa temperatura:
\begin{equation}
    \lambda_{\mathrm{max}}=\frac{b}{T}\quad \Rightarrow\quad E=h\nu=h\frac{c}{\lambda}=\frac{h c T}{b}=1.15\cdot10^{-3}\,\mathrm{eV}
\end{equation}
Successivamente l'energia disponibile nel centro di massa si calcola:
\begin{equation}
    S=\left(\mathbf{P}_p+\mathbf{P}_\gamma\right)^2=m_p^2+0+2E_pE_\gamma-2|\mathbf{p}_p|E_\gamma\cos{\theta}
\end{equation}
Per calcolare l'energia di soglia si calcola il caso in cui $S$ è massimo e quindi:
\begin{equation}
    S_\mathrm{max}: \quad \cos{\theta}=-1
\end{equation}
Di conseguenza l'espressione diventa:
\begin{equation}
    S_\mathrm{max}=m_p^2+2E_pE_\gamma+2|\mathbf{p}_p|E_\gamma
\end{equation}
Ora è possibile porre il protone come ultrarelativistico per semplificare i conti. Successivamente dopo aver controllato il risultato è possibile verificare se questo è vero oppure no. Quindi $|\mathbf{p}_p|\sim E_p$ ($E_p>>m_p$):
\begin{equation}
    S_\mathrm{max}=m_p^2+2E_pE_\gamma+2\sqrt{E_p^2-m_p^2}E_\gamma=m_p^2+4E_pE_\gamma=m^2_\Delta
\end{equation}
Sapendo che la massa della $\Delta$ è: $m_\Delta=1232\,\mathrm{MeV}$ posso calcolare l'energia di soglia del protone per indurre la reazione richiesta:
\begin{equation}
    E_p^{th}=\frac{m_\Delta^2-m_p^2}{4E_\gamma}=1.3\cdot10^{20}\,\mathrm{eV}
\end{equation}
Quindi supporre il protone come una particella ultrarelativistica è corretto.